%! Equivalent Representations of Polynomial Expressions — Unit 2, Lesson 2.6 — Warm-Up
\documentclass[10pt]{article}
\usepackage{precalculus-article}
\usepackage{precalculus-boxes}
\newcommand{\TallMath}[1]{$\displaystyle #1\rule[-1.4em]{0pt}{3.2em}$}

\begin{document}

\pageheader{Unit 2, Lesson 2.6}{Warm-Up}

\vspace{0.12in}

\begin{enumerate}[itemsep=12pt]

\item From Lesson 2.5: $p(x) = -2x^5 + 3x^2 - 7$.

\begin{enumerate}[label=(\alph*), itemsep=6pt]
  \item Degree: \blank{1.4cm} \qquad Leading coefficient: \blank{1.6cm}
  \item Describe the end behavior in limit notation.

        \vspace{0.05in}
        $\lim_{x \to \infty} p(x) = $ \blank{2cm} \qquad
        $\lim_{x \to -\infty} p(x) = $ \blank{2cm}
\end{enumerate}

\item From Lesson 2.3: $q(x) = (x-4)(x+2)^2$.

\begin{enumerate}[label=(\alph*), itemsep=6pt]
  \item List the real zeros of $q$ with the multiplicity of each.

        \vspace{0.05in}
        Zeros: \blank{6.5cm}
  \item What is the degree of $q$? \quad \blank{1.6cm}
\end{enumerate}

\item From Algebra 2: multiply and write the result in standard form.

\begin{enumerate}[label=(\alph*), itemsep=6pt]
  \item $(x+3)(x-5)$

        \begin{work}
          (x+3)(x-5) &= x^2 - 5x + 3x - 15 \\
                     &= x^2 - 2x - 15
        \end{work}

        $(x+3)(x-5) = $ \blank{4.5cm}
\end{enumerate}

\end{enumerate}

\end{document}
